\chapter{Another Language in a Zeno File}\label{ch-inline}

Some code is already written, and is in another language: a checksum
out of a board's firmware, a few lines of Python that use a library, a
function of JavaScript from a page. Since 0.9 an \verb|inline| block
holds it in a \verb|.zeno| file as it was written, and the program
calls it:

\begin{zeno}[file=ones.zeno, dir=blocks]
import c

let bits = inline c {
    #include <stdint.h>

    /* How many bits of x are 1. */
    int ones(uint32_t x) {
        int count = 0;
        for (; x != 0; x &= x - 1) {
            count++;
        }
        return count;
    }
}

print(bits.ones(0xA5))
\end{zeno}

What is between the first line and the last \verb|}| is C, and is not
read as Zeno. \verb|bits| is a library of C, as \verb|c.library| gives
one (\hyperref[ch-c-libraries]{Libraries of C}), and \verb|ones| is a
function of it.

There are three kinds of block:

\begin{center}
\begin{tabular}{llP{6.4cm}}
\toprule
Written & Holds & Gives \\
\midrule
\verb|inline c| & C & a library of C: the functions the block writes
  out \\
\verb|inline python| & Python & a module of Python: what the block
  defines \\
\verb|inline node| & JavaScript & what the block puts in
  \verb|exports| \\
\bottomrule
\end{tabular}
\end{center}

\begin{note}
Nothing changes for a program that has no block. A block is asked for
twice: by the word \verb|inline|, and by the \verb|import| of its
language at the top of the file. \verb|inline| is not a reserved
word, and \verb|let inline = 3| still means what it meant.
\end{note}

\section{Why Zeno is not a superset of C, or of Python}

It would be one way to use code that is already written: make every
program of C a program of Zeno, as every program of C is nearly one of
C++, and every program of JavaScript is one of TypeScript. Zeno does
not, for three reasons.

A superset has to take every program of the other language, for ever,
with every mistake that language allows. The first thing given up
would be the checks Zeno is for: a sum that cannot overflow unseen, a
name that cannot be \verb|none| unseen. A superset of two languages
cannot exist where they disagree, and \verb|a / b|, \verb|{}| and
\verb|'x'| each mean something else in C, Python and JavaScript. And a
superset of one language is a dialect of it: TypeScript cannot mend
the \verb|==| of JavaScript.

A block gets the same result for the person with the code. It is
used as it is, today, by that language's own compiler, and Zeno stays
one language with one set of rules.

\section{What it needs}

What the \verb|import| of the language needs, and nothing more: a C
compiler for \verb|inline c|, as for \verb|import c| and
\verb|zeno build|; Python 3 for \verb|inline python|; Node for
\verb|inline node|
(\hyperref[ch-other-languages]{Libraries of Other Languages}).

\section{Writing a block}

\begin{zenofrag}[parse=no]
let name = inline c {
    ...the other language...
}
\end{zenofrag}

\begin{itemize}
\item \textbf{A block is the whole value of a \texttt{let}.} What it
  has is then used through that name. It may be at the top of a file,
  in a function, or in a test.
\item \textbf{The first line ends with the \texttt{\{}.}
  \verb|inline|, the language, \verb|{|, and the end of the line.
\item \textbf{Everything inside is further in than the line that
  opens the block.} How far is up to you; \verb|zeno format| sets it
  one step in.
\item \textbf{The block ends at a line that holds only \texttt{\}}},
  as far in as the line that opened it.
\end{itemize}

That last rule is why a block needs no care. A \verb|}| of C is
always further in than the block's own, so it ends nothing. Zeno does
not look for the other language's strings and comments, and so cannot
mistake one:

\begin{zeno}[file=odd.zeno, dir=blocks]
import c

let odd = inline c {
    /* A } in a comment, */
    const char *closing(void) { return "} and one in a text"; }
}

print(odd.closing())
\end{zeno}

The indentation that every line of a block shares is not part of the
block. Each line keeps how far in it is from the others, which in
Python is part of the program.

\begin{why}
To find the end of a block by counting braces, Zeno would have to know
what a brace inside a string or a comment looks like in each language:
\verb|"}"|, \verb|'}'|, \verb|/* } */|, a template of JavaScript, a
regular expression, a raw string of Python. That is a small reader for
each language, each wrong somewhere, and the mistake it would make is
the worst kind: the block ends early, and the rest of the other
language's code is read as Zeno. A line that holds only \verb|}| at
the first line's own indentation is found with no knowledge of the
other language, and can be seen with the eye. It is the one place in
Zeno where indentation means something, and a line that breaks the
rule is refused, never guessed at. (Decision 61.3.)
\end{why}

\section{A block of C}

Each function that is written out in the block, with its body, and is
not \verb|static|, is a function of the library. Its types are read
from its own first line, so nothing is declared a second time:

\begin{zeno}[file=frame.zeno, dir=blocks]
import c

let frame = inline c {
    #include <stdint.h>
    #include <stddef.h>

    #define POLYNOMIAL 0xA001

    static uint16_t step(uint16_t crc) {
        return (crc & 1) ? (crc >> 1) ^ POLYNOMIAL : crc >> 1;
    }

    uint16_t crc16(const uint8_t *data, size_t length) {
        uint16_t crc = 0xFFFF;
        for (size_t i = 0; i < length; i++) {
            crc ^= data[i];
            for (int bit = 0; bit < 8; bit++) {
                crc = step(crc);
            }
        }
        return crc;
    }
}

let request: list[u8] = [0x01, 0x03, 0x00, 0x6B, 0x00, 0x03]
print(frame.crc16(request, len(request)))
\end{zeno}

\verb|crc16| is a function of \verb|frame|. \verb|step| is
\verb|static|, so it is the block's own, and so are
\verb|POLYNOMIAL|, a variable, a \verb|typedef| and a \verb|struct|:
they are there for the block's functions, and Zeno leaves them alone.

A call is checked before the program runs, as a call of a function of
Zeno is:

\begin{zeno}[file=calls.zeno, dir=blocks, cmd=zeno check {file}, fails=yes]
import c

let frame = inline c {
    #include <stdint.h>
    #include <stddef.h>

    static uint16_t step(uint16_t crc) { return crc >> 1; }

    uint16_t crc16(const uint8_t *data, size_t length) {
        uint16_t crc = 0xFFFF;
        for (size_t i = 0; i < length; i++) {
            crc = step(crc ^ data[i]);
        }
        return crc;
    }
}

print(frame.crc16(3, 1))
print(frame.crc16([1, 2]))
print(frame.step(1))
\end{zeno}

Values go across as they do for \verb|c.library|: a whole number is
checked to fit the type the function takes, a list is lent for a
pointer and copied back unless the pointer is to \verb|const|, text
goes as a \verb|const char *|, and any other pointer C gives back is
kept as it is. It is one table for both, and it is in the last
chapter. A header is included as in C: one of the system in
\verb|< >|, one of the program in \verb|" "|, found from the folder of
the \verb|.zeno| file.

A mistake in the block is the C compiler's to find, and it is told
where each line of the block was written, so what it says names the
\verb|.zeno| file:

\begin{zeno}[file=sensor.zeno, dir=blocks, fails=yes]
import c

let lib = inline c {
    int twice(int x) { return 2 * y; }
}

print(lib.twice(21))
\end{zeno}

The line is the line of the \verb|.zeno| file. The column is counted
from where the lines of the block begin, since the indentation they
share is not part of the block.

A block that does not compile stops the program where the block is.
No \verb|try| catches it: the program is wrong, and running it again
will not help.

\begin{center}
\begin{tabular}{P{6.2cm}P{6.6cm}}
\toprule
A block of C cannot have & What to do \\
\midrule
a function Zeno cannot call that is not \verb|static|: one that takes
  \verb|...|, a function, or a whole struct & mark it \verb|static| if
  only the block calls it, or write a function beside it that Zeno can
  call \\
no function Zeno can call & write one out, with its body, without
  \verb|static| \\
C++ & functions of C++ marked \verb|extern "C"|, built into a library,
  go to \verb|c.library| \\
\bottomrule
\end{tabular}
\end{center}

\section{What zeno build does with a block}

A build that writes C writes the block into the program's C, as it
was written. There is no library beside the program and nothing
between a call and the function: \verb|frame.crc16(...)| is a call of
\verb|crc16| in the same file, which the C compiler may put in line.

\begin{zeno}[file=frame.zeno, dir=blocks, show=no, cmd=zeno build {file} ;; ./frame ;; zeno build --no-collector --minimal {file} ;; ./frame ;; grep -c "^#line" frame.c]
import c

let frame = inline c {
    #include <stdint.h>
    #include <stddef.h>

    #define POLYNOMIAL 0xA001

    static uint16_t step(uint16_t crc) {
        return (crc & 1) ? (crc >> 1) ^ POLYNOMIAL : crc >> 1;
    }

    uint16_t crc16(const uint8_t *data, size_t length) {
        uint16_t crc = 0xFFFF;
        for (size_t i = 0; i < length; i++) {
            crc ^= data[i];
            for (int bit = 0; bit < 8; bit++) {
                crc = step(crc);
            }
        }
        return crc;
    }
}

let request: list[u8] = [0x01, 0x03, 0x00, 0x6B, 0x00, 0x03]
print(frame.crc16(request, len(request)))
\end{zeno}

Every build that writes C does it: a plain one, \verb|--firmware|,
\verb|--no-collector|, \verb|--owned|, \verb|--minimal|,
\verb|--release|, \verb|--fast|, \verb|--board|, and
\verb|zeno test --build| (\hyperref[ch-compiling]{Compiling}). The
last line above counts the lines of the built C that begin with
\verb|#line|. There are two: one before the block, which names the
\verb|.zeno| file and the line the block began on, and one after it,
which gives the lines that follow back to the C file. So a mistake
the C compiler finds, and a line a debugger stops at, is a line of the
\verb|.zeno| file.

Under \verb|zeno run| the block is built once into a small library,
kept, and called through the helper that \verb|import c| uses; it is
built again when the block changes.

\section{Assembly}

Assembly is written inside a block of C, as the C compiler takes it:

\begin{zenofrag}[parse=no]
import c

let cpu = inline c {
    #include <stdint.h>

    /* Interrupts off, on a Cortex-M. What PRIMASK was is given back. */
    uint32_t interrupts_off(void) {
        uint32_t was;
        __asm__ volatile ("mrs %0, primask\n cpsid i" : "=r" (was));
        return was;
    }
}
\end{zenofrag}

There is no \verb|inline asm|. Assembly has no types, so nothing in it
says what a routine takes or gives, and the C function around it is
what says so. It is also not one language: each processor has its own
instructions and each compiler its own way of writing them, and the C
compiler already knows both.

A whole file of assembly is given to \verb|c.library|, beside the
\verb|.c| files it takes:

\begin{zenofrag}
let start = c.library("startup.s", "void reset(void);")
\end{zenofrag}

The C compiler assembles it, a \verb|.s| as it is and a \verb|.S|
after reading its \verb|#include| and \verb|#define| lines.

\section{A block of Python}

\begin{zeno}[file=stats.zeno, dir=blocks]
import python

let stats = inline python {
    import statistics

    def summary(readings):
        """The middle and the spread of some readings."""
        return [statistics.median(readings), statistics.pstdev(readings)]
}

print(stats.summary([21.5, 22.0, 23.5]))
\end{zeno}

The block is a module of Python: what \verb|python.module| would give
if the lines were in a \verb|.py| file. Everything it defines is
reached with a dot: a function, a class, a variable, and a module the
block imported (\verb|stats.statistics|). Values go across as for
\verb|import python|.

The block runs when the program gets to it, once, top to bottom, as a
module does when it is imported. An error of Python in it stops the
program, and names the line of the \verb|.zeno| file:

\begin{zeno}[file=wrong.zeno, dir=blocks, fails=yes]
import python

let stats = inline python {
    def summary(readings:
        return max(readings)
}

print(stats.summary([1, 2]))
\end{zeno}

\section{A block of JavaScript}

\begin{zeno}[file=written.zeno, dir=blocks]
import node

let written = inline node {
    exports.number = (value, place) => new Intl.NumberFormat(place).format(value)
}

print(written.number(1234567.5, "en-IN"))
print(written.number(1234567.5, "de-DE"))
\end{zeno}

The block is a module of Node, as a \verb|.js| file is: what it puts
in \verb|exports|, or gives \verb|module.exports|, is what the block
gives, and \verb|require| finds the packages \verb|npm install| put in
the program's folder. It is a module of the older kind, with
\verb|require| and \verb|exports|; \verb|import| and \verb|export| are
not read there. Values go across as for \verb|import node|.

\section{What each command does with a block}

\begin{center}
\begin{tabular}{P{4.2cm}P{4.2cm}P{4.2cm}}
\toprule
 & \verb|inline c| & \verb|inline python|, \verb|inline node| \\
\midrule
\verb|zeno run|, \verb|zeno test| & built once, kept, and called
  through the C helper & given to the Python or the Node that the
  \verb|import| starts \\
\verb|zeno check| & the functions and every call of them are
  checked & checked when the program runs \\
\verb|zeno build|, and every build that writes C & written into the
  program's C & refused \\
\verb|zeno check --strict| & as \verb|zeno check| & refused: what it
  gives is not known before the program runs \\
\verb|zeno build --llvm| & refused & refused \\
\verb|zeno build --target| & refused & refused \\
the playground & refused: a page has no C compiler & refused \\
\verb|zeno format| & the lines of the block are left as they are, one
  step in & the same \\
\bottomrule
\end{tabular}
\end{center}

A refusal is an error with the line of the block, when the program is
checked or built. A block is never skipped.

\section{What Zeno cannot check}

The code in a block is the other language's, and Zeno checks none of
it. A pointer that is wrong in a block of C damages memory as it does
in C: a built program stops or goes wrong, and under \verb|zeno run|
the helper stops and Zeno says which call it was waiting for. The
blocks of a program are the lines to read then, and
\verb|grep -n "inline c"| finds them all, as \verb|grep -n unsafe|
finds the other lines of that kind (\hyperref[ch-unsafe]{Unsafe Code}).

\section{What is not there yet}

\begin{itemize}
\item \textbf{A library for a block.} A block of C that needs a
  library other than C's own (\verb|-lz|) has nowhere to say so. Until
  it has, the code goes in a \verb|.c| file given to \verb|c.library|
  with the library.
\item \textbf{A constant of a block.} A \verb|#define| or a variable
  of a block is not seen by Zeno; a function that gives it is.
  \verb|c.library| can declare a value.
\item \textbf{Small types given back as themselves.} A
  \verb|uint16_t| comes back as an \verb|int|, as for
  \verb|c.library|.
\item \textbf{Blocks of Ruby, PHP, Lua and C\#.} Each has an
  \verb|import| already and can be given a block in the same way.
\item \textbf{\texttt{zeno build --llvm}} with a block of C.
\end{itemize}

\section{How it is built}

The lexer decides: \verb|inline|, a name, a \verb|{| and the end of
the line begin a block, and it is the one place the lexer looks
further than two characters ahead. It reads the lines to the closing
one and makes of them one token, which holds the language and the
text. The parser allows that token only as the value of a
\verb|let|.

For a block of C, the checker finds the functions: it goes through
the text once, stepping over comments, strings, lines that begin with
\verb|#| and bodies between braces, and keeps the first line of each
function that is not \verb|static|. Those lines are what
\verb|c.library| would have been given, and from there a block is a
library: the checker, \verb|zeno run| and \verb|zeno build| use what
\verb|import c| built. A block of Python is compiled by Python's own
\verb|compile|, with as many empty lines before it as the
\verb|.zeno| file has above it, so that the line numbers Python
counts are the file's. The reasons are chapter 61 of
\verb|docs/decisions.md|, and \verb|examples/bits.zeno| and
\verb|examples/report.zeno| are two whole programs.

\section{Things to try}

\begin{enumerate}
\item Run \verb|examples/bits.zeno|, then build it and run what was
  built. Open \verb|bits.c| and find the block. What is on the line
  before it, and on the line after?
\item In the first program of this chapter, write \verb|count+++;| in
  the block and run it. Which file and line does the C compiler name?
\item Mark \verb|ones| as \verb|static| and run \verb|zeno check|.
  Read the message. Then write a second function that is not
  \verb|static| and calls it.
\item Write \verb|let m = inline python {| in a file with no
  \verb|import python|, and read what Zeno says.
\item Write a block of Python with a class in it, make an object of
  the class from Zeno, and call a method. Then do the same with a
  class of JavaScript and \verb|node.new|.
\item Put a \verb|}| on a line by itself in the middle of a block of
  C, as far in as the rest of the block. Does it end the block? Move
  it out to where the \verb|let| begins and run \verb|zeno check|
  again.
\end{enumerate}
